% Canonical Paper I source.
The affine cocycle retains the effect of order on evaluation.  We now place the
same information over a commutative charge plane.  This construction has two
distinct uses.  For a positive add--scale history, it records the total additive and
logarithmic multiplicative charges, and it supplies a weighted integral whose value restores
the translation part that the charge endpoint alone forgets.  The distinction
between these two roles is essential: the accumulative commutative space is a
shadow of the presented history, not another presentation of its full arithmetic
syntax.

Throughout this section an elementary additive step and an elementary scaling
step are denoted by
\[
    \mathsf A_p(x)=x+p,
    \qquad
    \mathsf M_q(x)=e^q x,
    \qquad p,q\in\mathbb R.
\]
For a history \(\gamma=(g_1,\ldots,g_n)\), the entries are performed from left
to right, so that
\[
    \nu_x(\gamma)=g_n\circ\cdots\circ g_1(x).
\]
Negative values of \(p\) and \(q\) are allowed.  Thus the same notation covers
inverse translations and inverse positive scalings.
In this section an \emph{add--scale history} means a finite word in these two
families of generators.  Such a word evaluates in
\(\operatorname{Aff}^+(1,\mathbb R)\).  The ACS construction is attached to
this presented word; it is not asserted for negative multipliers or for an
arbitrary affine syntax before that syntax has been written in add--scale form.

\subsection{The charge path}

\begin{definition}[Accumulative commutative space and charge path]
\label{def:acs}
The \emph{accumulative commutative space}, abbreviated ACS, is the oriented
plane
\[
    \mathcal A=\mathbb R^2_{(A,M)},
    \qquad \text{with orientation } dA\wedge dM.
\]
An additive step \(\mathsf A_p\) has charge increment \((p,0)\), and a scaling
step \(\mathsf M_q\) has charge increment \((0,q)\).  Starting at
\((A_0,M_0)=(0,0)\), accumulate these increments in chronological order.  The
resulting oriented broken path is denoted by \(C_\gamma\), and its terminal
point by
\[
    (A_\gamma,M_\gamma)
      =\left(\sum_{g_i=\mathsf A_{p_i}}p_i,
              \sum_{g_i=\mathsf M_{q_i}}q_i\right).
\]
\end{definition}

The terminal charge is insensitive to order, whereas the broken path generally
is not.  In particular, two histories may have the same ACS endpoint while
inducing different affine maps.  The missing information is recovered by the
following integral formula.

\begin{proposition}[Direct ACS evaluation]
\label{prop:acs-evaluation}
For every add--scale history \(\gamma\) and every initial value \(x\),
\begin{equation}
    \nu_x(\gamma)
      =e^{M_\gamma}
        \left(x+\int_{C_\gamma}e^{-M}\,dA\right).
    \label{eq:direct-acs-evaluation}
\end{equation}
\end{proposition}

\begin{proof}
Let an additive step \(\mathsf A_{p_i}\) occur when the cumulative
multiplicative charge is \(M_{i-1}\).  Along the corresponding horizontal edge
of \(C_\gamma\),
\[
    \int e^{-M}\,dA=p_i e^{-M_{i-1}}.
\]
After that addition, the remaining multiplicative steps have total charge
\(M_\gamma-M_{i-1}\).  Consequently the contribution of \(p_i\) to the final
value is
\[
    p_i e^{M_\gamma-M_{i-1}}
      =e^{M_\gamma}p_i e^{-M_{i-1}}.
\]
The initial value contributes \(e^{M_\gamma}x\), and vertical ACS edges
contribute nothing to the \(dA\)-integral.  Summing over the additive steps
gives \eqref{eq:direct-acs-evaluation}.  The argument uses oriented edge
integrals, so it also covers negative increments.
\end{proof}

The kernel \(e^{-M}\) is therefore forced by the affine cocycle: it normalizes
an additive contribution to the source frame, before the common terminal scale
\(e^{M_\gamma}\) returns it to the target frame.

\subsection{Compatibility and relative torsion}

There are two useful compatibility conditions, and they should not be
conflated.  Histories \(\gamma\) and \(\delta\) are \emph{scale-compatible} if
\[
    M_\gamma=M_\delta.
\]
They are \emph{charge-compatible} if, in addition,
\[
    A_\gamma=A_\delta.
\]
Scale compatibility makes their affine maps have the same linear part.  Charge
compatibility makes their ACS paths have both endpoints in common and hence
makes \(C_\gamma-C_\delta\) a closed oriented one-chain.

\begin{definition}[Relative arithmetic torsion]
\label{def:relative-torsion}
For scale-compatible add--scale histories \(\gamma\) and \(\delta\), their \emph{relative
arithmetic torsion} is
\begin{equation}
    \tau(\gamma,\delta)
       :=\xi_\gamma-\xi_\delta,
    \label{eq:relative-torsion}
\end{equation}
where \(\xi_\gamma\) and \(\xi_\delta\) are the target-frame translation
terms of the two affine operators.  This relative translation defect is the
primary torsion invariant used below.
\end{definition}

\begin{proposition}
\label{prop:relative-torsion-independent}
For every initial value \(x\), relative arithmetic torsion equals the endpoint
difference
\[
    \tau(\gamma,\delta)
      =\nu_x(\gamma)-\nu_x(\delta),
\]
and is therefore independent of \(x\).
If \(M_\gamma=M_\delta=M_*\), then
\begin{equation}
    \tau(\gamma,\delta)
      =e^{M_*}
        \left(
          \int_{C_\gamma}e^{-M}\,dA
          -\int_{C_\delta}e^{-M}\,dA
        \right).
    \label{eq:relative-torsion-boundary-preparation}
\end{equation}
\end{proposition}

\begin{proof}
Apply Proposition~\ref{prop:acs-evaluation} to both histories.  Their initial
terms are respectively \(e^{M_\gamma}x\) and \(e^{M_\delta}x\), which cancel
under scale compatibility.  The remaining terms give
\eqref{eq:relative-torsion-boundary-preparation}.
\end{proof}

Temporal reversal gives one possible comparison, but it is not built into the
definition.  The pair \((\gamma,\delta)\) may consist of any two
scale-compatible add--scale histories.  This extension is important because order defects
occur between histories that are neither reversals nor planar mirrors of one
another.

\subsection{The torsion--Stokes formula}

Suppose now that \(\gamma\) and \(\delta\) are charge-compatible, with common
terminal charge \((A_*,M_*)\).  Define the target-normalized one-form
\begin{equation}
    \eta_*:=e^{M_*-M}\,dA
    \label{eq:target-normalized-acs-form}
\end{equation}
on the ACS.  Relative torsion is then the integral of \(\eta_*\) over the
closed one-chain \(C_\gamma-C_\delta\).  The chosen order of the two paths
fixes the sign.

\begin{theorem}[Weighted torsion--Stokes theorem]
\label{thm:torsion-stokes}
Let \(\gamma\) and \(\delta\) be charge-compatible add--scale histories with
common terminal multiplicative charge \(M_*\).  Let
\(\Sigma_{\gamma,\delta}\) be any oriented piecewise smooth singular
two-chain in the ACS satisfying
\[
    \partial\Sigma_{\gamma,\delta}=C_\gamma-C_\delta.
\]
Then
\begin{align}
    \tau(\gamma,\delta)
      &=\int_{C_\gamma-C_\delta}\eta_*
        \label{eq:torsion-boundary-integral}\\
      &=\int_{\Sigma_{\gamma,\delta}}
          e^{M_*-M}\,dA\wedge dM.
        \label{eq:torsion-area-integral}
\end{align}
The value of the area integral is independent of the chosen filling chain.
\end{theorem}

\begin{proof}
Equation \eqref{eq:relative-torsion-boundary-preparation} and charge
compatibility give
\[
    \tau(\gamma,\delta)
       =\int_{C_\gamma-C_\delta}e^{M_*-M}\,dA
       =\int_{C_\gamma-C_\delta}\eta_*.
\]
With the orientation \(dA\wedge dM\), direct differentiation yields
\[
    d\eta_*
      =d(e^{M_*-M})\wedge dA
      =-e^{M_*-M}dM\wedge dA
      =e^{M_*-M}dA\wedge dM.
\]
Stokes' theorem for piecewise smooth chains therefore gives
\[
    \int_{C_\gamma-C_\delta}\eta_*
      =\int_{\partial\Sigma_{\gamma,\delta}}\eta_*
      =\int_{\Sigma_{\gamma,\delta}}d\eta_*,
\]
which is \eqref{eq:torsion-area-integral}.  Such a filling exists because the
ACS plane is contractible.  Finally, if \(\Sigma\) and \(\Sigma'\) have the
same boundary, applying Stokes' theorem to each shows
\[
    \int_\Sigma d\eta_*
       =\int_{\partial\Sigma}\eta_*
       =\int_{\partial\Sigma'}\eta_*
       =\int_{\Sigma'}d\eta_*.
\]
This also proves filling independence.  The chain formulation remains valid
when the broken paths self-intersect or carry signed increments; no informal
notion of an enclosed region is required.
\end{proof}

Charge compatibility is needed for the closed-chain and area formulas, but not
for Definition~\ref{def:relative-torsion}.  A merely scale-compatible pair
still has a well-defined translation defect, even when the endpoints of its ACS
paths have different \(A\)-coordinates.

Figure~\ref{fig:acs-compatible-paths-weighted-filling} depicts the chain
orientation used in Theorem~\ref{thm:torsion-stokes}.  It is a schematic of
the theorem's data rather than an argument for filling independence.

\begin{figure}[t]
\centering
\begin{tikzpicture}[
  x=1.15cm,y=1.05cm,
  axis/.style={-{Latex[length=2mm]},thick,black!70},
  gamma/.style={-{Latex[length=2mm]},very thick,red!65!black},
  delta/.style={-{Latex[length=2mm]},very thick,blue!65!black},
  every node/.style={font=\small}
]
  \coordinate (O) at (0,0);
  \coordinate (g1) at (2.2,0);
  \coordinate (g2) at (2.2,0.9);
  \coordinate (g3) at (4,0.9);
  \coordinate (P) at (4,2.7);
  \coordinate (d1) at (0,1.5);
  \coordinate (d2) at (1.3,1.5);
  \coordinate (d3) at (1.3,2.7);

  \path[fill=violet!11,draw=none]
    (O)--(g1)--(g2)--(g3)--(P)--(d3)--(d2)--(d1)--cycle;
  \draw[axis] (-0.25,0) -- (4.75,0) node[right] {$A$};
  \draw[axis] (0,-0.25) -- (0,3.25) node[above] {$M$};

  \draw[gamma] (O)--(g1);
  \draw[gamma] (g1)--(g2);
  \draw[gamma] (g2)--(g3);
  \draw[gamma] (g3)--(P);
  \draw[delta] (O)--(d1);
  \draw[delta] (d1)--(d2);
  \draw[delta] (d2)--(d3);
  \draw[delta] (d3)--(P);

  \fill (O) circle (1.7pt) node[below left] {$(0,0)$};
  \fill (P) circle (1.7pt) node[above right] {$(A_*,M_*)$};
  \node[red!65!black,fill=white,inner sep=1pt] at (3.3,0.72) {$C_\gamma$};
  \node[blue!65!black,fill=white,inner sep=1pt] at (0.72,1.75) {$C_\delta$};
  \node[align=center,violet!55!black] at (2.55,1.8)
    {$\Sigma_{\gamma,\delta}$\\
     $e^{M_*-M}\,dA\wedge dM$};
  \node[font=\scriptsize,anchor=west] at (4.35,1.72)
    {$\partial\Sigma=C_\gamma-C_\delta$};

  \draw[-{Latex[length=1.7mm]},black!55] (4.55,0.35)--(4.95,0.35)
    node[right,font=\scriptsize] {$dA$};
  \draw[-{Latex[length=1.7mm]},black!55] (4.55,0.35)--(4.55,0.8)
    node[above,font=\scriptsize] {$dM$};
\end{tikzpicture}
\caption{Two charge-compatible ACS histories.  Both paths run from $(0,0)$
to the common charge endpoint $(A_*,M_*)$.  Traversing $C_\gamma$ forward and
$C_\delta$ backward gives the positive boundary of the displayed filling for
the orientation $dA\wedge dM$; the density is weighted by $e^{M_*-M}$.}
\label{fig:acs-compatible-paths-weighted-filling}
\end{figure}

\subsection{Explicit order comparisons}

\paragraph{The elementary rectangle.}
Let
\[
    \gamma=(\mathsf A_p,\mathsf M_q),
    \qquad
    \delta=(\mathsf M_q,\mathsf A_p).
\]
The path \(C_\gamma-C_\delta\) is the positively oriented boundary of the
rectangle \([0,p]\times[0,q]\) when \(p,q>0\).  Direct evaluation and the
weighted-area formula agree:
\begin{align*}
    \tau(\gamma,\delta)
      &=e^q(x+p)-(e^q x+p)\\
      &=p(e^q-1)\\
      &=\int_0^p\int_0^q e^{q-M}\,dM\,dA.
\end{align*}
Thus ``add, then multiply'' minus ``multiply, then add'' is positive under the
orientation fixed in Definition~\ref{def:acs}.  In the flow notation used in
the next chapter, \(p=\mu h\) and \(q=\lambda k\), so the defect is
\[
    \mu h(e^{\lambda k}-1).
\]

\paragraph{A four-step staircase.}
For
\[
    \gamma=(\mathsf A_p,\mathsf M_r,\mathsf A_q,\mathsf M_s),
    \qquad
    \delta=(\mathsf M_r,\mathsf A_p,\mathsf M_s,\mathsf A_q),
\]
the histories have the same total charge \((p+q,r+s)\), but \(\delta\) is not
the temporal reversal of \(\gamma\).  Their evaluations are
\[
    \nu_x(\gamma)=e^{r+s}x+p e^{r+s}+q e^s,
\]
\[
    \nu_x(\delta)=e^{r+s}x+p e^s+q.
\]
Consequently
\begin{equation}
    \tau(\gamma,\delta)
      =p e^s(e^r-1)+q(e^s-1).
    \label{eq:four-step-staircase-torsion}
\end{equation}
For positive \(p,q,r,s\), both terms are positive weighted strips in an ACS
filling.  Formula \eqref{eq:four-step-staircase-torsion} also shows how a later
scale changes the weight of an earlier commutation defect.

\paragraph{Equal charges without reversal.}
The pair
\[
    (\mathsf A_p,\mathsf M_r,\mathsf A_q),
    \qquad
    (\mathsf A_q,\mathsf M_r,\mathsf A_p)
\]
has common charge \((p+q,r)\), and its relative torsion is
\[
    (p-q)(e^r-1).
\]
It may be positive, negative, or zero.  In particular, equal charge does not
determine torsion, while zero torsion does not imply equality of marked
histories.

\subsection{What the ACS forgets}

At the formal level, the map
\[
    \gamma\longmapsto(A_\gamma,M_\gamma)
\]
is an order-forgetting charge homomorphism from concatenated marked histories
to the additive group \(\mathbb R^2\).  It is therefore analogous to an
abelianization map.  This statement concerns the formal charge language only.
It does not identify the symbolic history space with the evaluated affine
group, and it does not assert that either object is a free group without an
additional choice of formal generators and relations.

The content of Theorem~\ref{thm:torsion-stokes} is precisely what survives this
loss of order: the endpoint charges specify a common frame, while the weighted
filling measures the translation defect between two histories.  The next
chapter retains the current value as a third coordinate and identifies the
infinitesimal form of the same defect with contact curvature.
